\documentclass[10pt,a4paper]{amsart}
\usepackage[margin=19mm]{geometry}
\usepackage{amsmath,amssymb,mathtools}
\usepackage[scr=rsfs,cal=euler]{mathalfa}
\usepackage{xfrac}
\usepackage[unicode,hidelinks,bookmarks=false]{hyperref}
\newcommand{\ie}{i.e.\ }
\newcommand{\cf}{cf.\ }
\newcommand{\resp}{respectively\ }
\newcommand{\Subsection}{\subsection}
\providecommand{\nobreakdash}{\nobreak\hbox{-}}
\setlength{\emergencystretch}{2em}
\multlinegap0pt
\usepackage{bm}
\usepackage{bbm}
\usepackage{dsfont}
\usepackage{xspace}
\usepackage{cancel}
\usepackage{mathtools}
\usepackage{amsthm}
\makeatletter
\@ifundefined{theorem}{\newtheorem{theorem}{Theorem}}{}
\makeatother
\title[A Spinorial Heat Flow Framework for Geometric Degeneration o]{A Spinorial Heat Flow Framework for Geometric Degeneration on $3$-Manifolds}
\author{Ferhat Ta\c{s}}
\date{Source: arXiv:2512.19308v2 (CC BY 4.0)}
\begin{document}
\maketitle

\noindent\emph{Source and licence.} A Spinorial Heat Flow Framework for Geometric Degeneration on $3$-Manifolds. Version arXiv:2512.19308v2 (CC BY 4.0), distributed under the Creative Commons Attribution 4.0 International licence, as recorded in the arXiv licence metadata for this exact version and shown on the abstract page https://arxiv.org/abs/2512.19308v2. Full rights notice: licenses/arxiv-2512.19308-CC-BY-4.0.txt.\par\smallskip

\noindent\textit{Editorial scope.} Counted unit: the Theorem whose source label is \texttt{None} in the article (source0.tex lines 426--442, source label \texttt{None}), delivered together with the article's own complete original proof of it (source0.tex lines 444--488, one \texttt{proof} environment of 45 lines). No mathematics is written, repaired or re-derived: statement and proof are the article's own text line for line. Required context retained uncounted: none. Every definition, display and label of those lines is retained unchanged. Omitted from this package and not redistributed: 1--425, 443--443, 489--854. Nothing inside a retained excerpt is omitted or altered: every retained body is delivered exactly as the article's lines are written, so no omission receipt of any kind accompanies this package. Citations: the retained excerpts carry no citation command at all, so no bibliography entry of the article is transcribed and no reference is redistributed. Reference adaptation: none was needed and none was made; every reference inside the delivered unit resolves inside it, so no unresolved reference or citation is produced. Printed slips kept literal: no printed word, symbol, hypothesis or number of the counted statement or of its proof was altered. Layout and class: the article's own class is \texttt{amsart} with packages including inputenc, fontenc, babel, amsmath, amssymb, amsfonts, amsthm, physics, mathtools, geometry, graphicx, xcolor, tikz, pgfplots; the wrapper is the lane's pinned \texttt{amsart} editorial wrapper at 10pt on A4, which restates the theorem-like environments the retained text needs and adds the article's own macro definitions that the retained text needs (none), with extra wrapper packages (bm, bbm, dsfont, xspace, cancel, mathtools). The delivered numbering is the unit's own. Assets: no retained span contains a figure environment, a graphics-inclusion command, a PDF-page inclusion, a table or a TikZ picture; no image file is copied into this package, the package declares no asset dependency and no image file is redistributed with it. Licence: version arXiv:2512.19308v2 of this article is distributed under the Creative Commons Attribution 4.0 International licence, as recorded in the arXiv licence metadata for this exact version and shown on the abstract page https://arxiv.org/abs/2512.19308v2. A full rights notice is delivered at licenses/arxiv-2512.19308-CC-BY-4.0.txt. Characters: every retained excerpt is pure ASCII, so the delivered engine, XeLaTeX with its default Latin Modern text font, reports no missing character.
\begin{theorem}[Short-time existence of weak solutions]
Let $(M^3,g_0)$ be a closed spin manifold and let
$\psi_0 \in H^1(\rho_0^\alpha\, dV_{g_0})$ be an initial spinor field,
where $\rho_0 := |\psi_0|$ and $\alpha>0$.
Then there exists a time $T>0$ and a weak solution
\[
\psi \in L^2(0,T; H^1(\rho^\alpha\, dV_{g_0})) \cap
H^1(0,T; H^{-1}(\rho^\alpha\, dV_{g_0}))
\]
to the spinorial heat flow
\[
\partial_t \psi = - D_{g(\psi)}^{\,2}\psi,
\]
in the sense of distributions.
Moreover, $\psi(t)$ depends continuously on the initial data in the
weighted $L^2$ topology.
\end{theorem}

\begin{proof}[Sketch of the proof]
The proof proceeds by a regularization and compactness argument.

\medskip
\noindent
\textbf{Step 1: Regularization.}
For $\varepsilon>0$, consider the uniformly parabolic regularized equation
\[
\partial_t \psi_\varepsilon
= -\big(D_{g(\psi_\varepsilon)}^{\,2} + \varepsilon \Delta_{g_0}\big)\psi_\varepsilon,
\]
with smooth initial data $\psi_{\varepsilon}(0)=\psi_0^\varepsilon$
approximating $\psi_0$ in $H^1(\rho_0^\alpha dV_{g_0})$.
Standard parabolic theory yields a unique smooth solution
$\psi_\varepsilon$ on a time interval $[0,T_\varepsilon)$.

\medskip
\noindent
\textbf{Step 2: Uniform weighted energy estimates.}
Using the conformal structure $g(\psi)=\rho^2 g_0$ and the Lichnerowicz formula,
one derives weighted energy inequalities of the form
\[
\frac{d}{dt}\int_M \rho^\alpha |\psi_\varepsilon|^2\, dV_{g_0}
+ \int_M \rho^\alpha |\nabla \psi_\varepsilon|^2\, dV_{g_0}
\le C \int_M \rho^\alpha |\psi_\varepsilon|^2\, dV_{g_0},
\]
where the constant $C$ is independent of $\varepsilon$.
These estimates imply uniform bounds in
$L^2(0,T; H^1(\rho^\alpha dV_{g_0}))$.

\medskip
\noindent
\textbf{Step 3: Compactness and limit passage.}
By weak compactness, there exists a subsequence $\psi_{\varepsilon_j}$
converging weakly to a limit $\psi$ in the weighted Sobolev spaces above.
Passing to the limit in the weak formulation yields a distributional
solution of the original degenerate equation.

\medskip
\noindent
\textbf{Step 4: Continuity and uniqueness.}
Continuity with respect to the initial data follows from the weighted
energy inequality.
Uniqueness is not claimed at this level of regularity.
\end{proof}

\end{document}
